\documentclass[11pt]{article}

\usepackage{series}
\usepackage{amsmath,amssymb,amsfonts}
\usepackage{geometry}
\usepackage{hyperref}
\usepackage{enumitem}
\usepackage{array}
\geometry{margin=1in}

\title{Modal Triplet Theory:\\
Admissibility, Encodings, and the Structure of Physical Description\\
\large A Typed and Tiered Corpus Roadmap}
\author{Peter Nero}
\date{July 2026}

\begin{document}
\maketitle

\begin{abstract}
This paper is the organizational roadmap for Modal Triplet Theory (MTT).  It
does not add a new physical derivation.  It fixes the types, dependency order,
and proof tiers needed to read the corpus without promoting analogies or
target-compatible reconstructions into source theorems.  A projection, a
representative section, an exact upper decoder, autonomous reduced evolution,
and effective-state merger solve different mathematical problems.  Ordinary
noninjectivity does not remove a right inverse; autonomous descent instead
requires preservation of projection fibers.

For physical ten-dimensional realizations we adopt the canonical convention
$\pi:M_{10}\to Y_4$ with six-dimensional fiber $X_6$.  The modal triplet is
represented by typed vertical bundles, operators, and projectors.  It is not
three extra coordinate factors.  A compact shared circle is line-bundle phase
or holonomy data and is not Lorentzian time or a seventh internal coordinate.
Circle-Lens-Nil is retained as a useful obstruction taxonomy, not an
exhaustive theorem forcing unique gravity, gauge, and quantum responses.  We
prove the corresponding conditional $4+6$ realization statement and state
the exact limits of that result.

The roadmap classifies relations to quantum mechanics, quantum field theory,
general relativity, the Standard Model, strings, quantum gravity, and AQFT by
what is typed, reconstructed, selected, calibrated, or still open.  Failure of
one controlled chart means that chart has lost its certified domain; it does
not imply that every mathematical or physical description has ceased to
exist.  Paper-level status is a dated snapshot, while the current A/B ledger
and selected source hashes remain authoritative.
\end{abstract}

\tableofcontents
\newpage

\section{Purpose, Authority, and Scope}

MTT studies reduced descriptions whose existence, continuation, and physical
interpretation are conditional.  The corpus contains foundational functional
analysis, finite and continuum realizations, reconstructions of established
formalisms, numerical packets, and interpretive proposals.  These objects do
not all have the same logical status.

This paper has four purposes:
\begin{enumerate}[label=\arabic*.]
\item define a common typed dictionary;
\item record the dependency order among structural, realization, and physical
claims;
\item classify the strength of relations to established physical frameworks;
\item provide a stable reading guide while delegating live status to the
curated A/B ledger.
\end{enumerate}

\begin{remark}[Authority rule]
A paper is evidence for the results it proves, but it is not automatically the
current status authority for later work.  When a selected successor,
corrigendum, or current A/B row exists, its source path and hash control
status.  Corpus search is evidence discovery, not theorem promotion.
\end{remark}

\subsection{Three distinct levels}

Throughout the roadmap we distinguish:
\begin{enumerate}[label=\textbf{T\arabic*.}]
\item \textbf{Structural level:} typed spaces, maps, domains, margins, and
conditional theorems about reduced description;
\item \textbf{Realization level:} a geometry, bundle, connection, operator, or
finite algebra implementing some structural data;
\item \textbf{Physical level:} a selected action or dynamics, state or measure,
observable map, uncertainties, and comparison with data.
\end{enumerate}
Success at T1 does not imply T2, and success at T2 does not imply T3.

\subsection{What this paper does not claim}

This roadmap does not derive spacetime, a probability law, the Born rule,
Einstein dynamics, the Standard Model values, a string vacuum, or a
nonperturbative quantum theory of gravity.  It records which parts of those
programs are typed or conditionally reconstructed and which source obligations
remain.

\section{Reduced Description: The Core Dictionary}

\subsection{Upper and reduced data}

Let $\mathcal X$ be an upper configuration space, $\mathcal Y$ a reduced
space, and
\[
\Pi:\mathcal X\longrightarrow\mathcal Y
\]
a declared reduction map.  Let $F:\mathcal X\to\mathcal X$ be an upper
evolution or update.

\begin{definition}[Representative section]
A representative section is a map $s:\mathcal Y\to\mathcal X$ satisfying
\[
\Pi\circ s=\operatorname{id}_{\mathcal Y}.
\]
It chooses one upper representative from each reduced state in its domain.
\end{definition}

\begin{definition}[Exact upper decoder]
An exact upper decoder is a map $D:\mathcal Y\to\mathcal X$ satisfying
\[
D\circ\Pi=\operatorname{id}_{\mathcal X}
\]
on its declared upper domain.  Such a decoder can exist only where $\Pi$ is
injective.
\end{definition}

\begin{definition}[Autonomous reduced evolution]
An autonomous reduced evolution is a map $G:\mathcal Y\to\mathcal Y$ such that
\[
\Pi\circ F=G\circ\Pi.
\]
\end{definition}

\begin{definition}[Effective merger]
Two upper states merge effectively when $x_1\ne x_2$ but
\[
\Pi(x_1)=\Pi(x_2).
\]
This is a statement about reduced distinguishability, not automatically about
microscopic destruction or probability.
\end{definition}

These four notions must not be interchanged.

\subsection{The exact descent criterion}

\begin{theorem}[Factor-through criterion]
\label{thm:descent}
There exists a map $G:\Pi(\mathcal X)\to\Pi(\mathcal X)$ satisfying
$\Pi F=G\Pi$ if and only if $F$ preserves the fibers of $\Pi$:
\[
\Pi(x_1)=\Pi(x_2)
\quad\Longrightarrow\quad
\Pi(Fx_1)=\Pi(Fx_2).
\]
\end{theorem}

\begin{proof}
If $G$ exists, applying $\Pi F=G\Pi$ to $x_1$ and $x_2$ proves fiber
preservation.  Conversely, define
\[
G(\Pi x):=\Pi(Fx).
\]
Fiber preservation makes this definition independent of the chosen
representative, and the factorization identity follows.
\end{proof}

\begin{proposition}[Noninjectivity does not remove a right inverse]
\label{prop:rightinverse}
A noninjective surjection may have a representative section.
\end{proposition}

\begin{proof}
The projection $\Pi:\mathbb R^2\to\mathbb R$,
$\Pi(x_1,x_2)=x_1$, is noninjective and has the smooth section
$s(y)=(y,0)$.  Therefore noninjectivity alone cannot prove absence of a right
inverse.  Regular, global, equivariant, or stable sections require their own
hypotheses.
\end{proof}

The corrected A0 program treats failure of descent, merger, failure of exact
decoding, and ill-conditioning of representative continuation separately
\cite{Nero2026_MTT_A0_StructuralCore}.

\subsection{Probability and irreversibility}

A many-to-one reduction does not by itself define probabilities.  To obtain a
reduced Markov kernel or Born-type rule one must supply a measure or state and
prove the relevant pushforward.  Likewise, effective merger does not by itself
select one ontic history, define entropy, or establish a physical arrow of
time.

\section{Admissibility and Controlled Boundaries}

\subsection{A declared control contract}

An admissible chart must name the properties it controls.  Typical margins
include:
\begin{itemize}
\item a spectral gap $\gamma>0$;
\item boundedness and regularity of a Riesz or coherent projector;
\item a descent defect
\[
\delta_{\rm desc}
=\sup_{\Pi x_1=\Pi x_2}
d_{\mathcal Y}\!\left(\Pi Fx_1,\Pi Fx_2\right);
\]
\item conditioning of a representative section;
\item invariance and completeness of a fixed-point basin;
\item coherent stability or a Lyapunov margin;
\item truncation, tail, or approximation error bounds.
\end{itemize}

\begin{definition}[Controlled encoding chart]
A controlled encoding chart is a tuple
\[
\mathfrak C=(U,\Pi,\mathcal O,F,\mathfrak m,\varepsilon)
\]
where $U\subseteq\mathcal X$ is the declared domain, $\mathcal O$ is the
observable family, $\mathfrak m$ is a list of named control margins, and
$\varepsilon$ is the associated error budget.
\end{definition}

\begin{definition}[Chart boundary]
The boundary of a controlled chart is the locus at which at least one declared
margin reaches its admissible threshold or its continuation theorem ceases to
apply.
\end{definition}

\begin{proposition}[Local loss of control is not universal nonexistence]
\label{prop:boundary}
Failure of one chart contract proves only that this chart is no longer
certified on the failed domain.  It does not prove that every other
mathematical or physical description is meaningless there.
\end{proposition}

\begin{proof}
The premises concern the objects and margins of one chart.  Another chart may
use a different projection, domain, observable algebra, or stability theorem.
Without a theorem excluding all such alternatives, universal nonexistence
does not follow.
\end{proof}

\subsection{Selection fronts}

A selection front is therefore a boundary-layer regime for a declared
encoding.  It may exhibit large condition numbers, closing gaps, metastability,
or loss of robust section continuation.  A transition, reset, detector
instrument, or post-boundary branch must be supplied separately.  Boundary
crossing alone does not choose an outcome.

\section{Canonical Geometric Convention}

\subsection{Physical ten-dimensional specialization}

When a physical ten-dimensional realization is used, the canonical convention
is a bundle or locally trivial fibration
\[
\pi:M_{10}\longrightarrow Y_4
\]
with a four-dimensional Lorentzian base $Y_4$ and a six-dimensional compact
Riemannian fiber $X_6$.  In a global product specialization,
\[
M_{10}=Y_4\times X_6.
\]
This is a selected realization convention, not a theorem that all MTT models
must be ten-dimensional.

\subsection{Coordinate, bundle, and operator types}

The following are distinct:
\begin{itemize}
\item coordinate factors of $X_6$;
\item vector or principal bundles over $X_6$;
\item a Hermitian shared line $L_{\rm shared}$ and its connection;
\item vertical operators on sections of declared bundles;
\item spectral projectors and their ranges;
\item finite carrier modules and matrix representations.
\end{itemize}
A line bundle has rank one in each fiber but adds no coordinate dimension.
An operator rank profile is not a manifold-dimension count.

\subsection{A conditional \texorpdfstring{$4+6$}{4+6} theorem}

\begin{proposition}[Conditional transverse-curvature realization]
\label{prop:fourplussix}
Assume:
\begin{enumerate}[label=(\roman*)]
\item a four-dimensional base $Y_4$;
\item a fiber
\[
X_6=F_1\times F_2\times F_3
\]
with each $F_i$ an oriented two-manifold;
\item a line bundle with connection on each $F_i$ whose curvature
$\Omega_i$ is a nonzero area form.
\end{enumerate}
Then $M=Y_4\times X_6$ is ten-dimensional, and the pulled-back curvature
forms $p_i^\ast\Omega_i$ are linearly independent.
\end{proposition}

\begin{proof}
Dimensions add:
\[
\dim M=4+2+2+2=10.
\]
If $\sum_i a_i p_i^\ast\Omega_i=0$, restrict the equality to tangent vectors
of $F_j$ while holding the other factors fixed.  Every term except the
$j$th vanishes, giving $a_j\Omega_j=0$.  Since $\Omega_j$ is nonzero,
$a_j=0$.  This holds for each $j$.
\end{proof}

\begin{remark}[Limit of Proposition~\ref{prop:fourplussix}]
The proposition proves a lawful ten-dimensional realization after its
four-dimensional base, three transverse surfaces, and three nonzero curvature
forms are assumed.  It does not prove that MTT uniquely selects those
assumptions, that ten dimensions are necessary, or that other realizations are
excluded.
\end{remark}

\subsection{The shared circle}

The shared circle is represented by common $U(1)$ phase or holonomy data,
typically through a Hermitian line bundle with connection.  It is counted once.
It is not identified with physical Lorentzian time.  A flat connection can
have nontrivial global holonomy, so circle response is not synonymous with
nonzero curvature.

\subsection{Local \texorpdfstring{$3\times3$}{3x3} and global
\texorpdfstring{$1+2+3$}{1+2+3} data}

A rank-three comparison field belongs to
\[
Q\in\Gamma(\operatorname{Hom}(TP,TI)).
\]
Its nine components decompose locally into three orientation directions and
six symmetric strain directions,
\[
\operatorname{Mat}(3,\mathbb R)
=\mathfrak{so}(3)\oplus\operatorname{Sym}(3,\mathbb R),
\qquad
6=1+2+3
\]
after a flag is selected.  This is a representation decomposition, not
manifold-dimension multiplication.  The selected q79 finite carrier also has
a $1<2<3$ rank profile, but rank agreement alone does not construct the
connection-, operator-, or Hessian-preserving continuum intertwiner
\cite{Nero2026_MTT_C_Realizations}.

\section{Circle-Lens-Nil: Taxonomy, Not Exhaustiveness}

\subsection{Recurring obstruction profiles}

Circle-Lens-Nil (CLN) is retained as a useful vocabulary:
\begin{itemize}
\item \textbf{Circle:} loop, recurrence, phase, or holonomy consistency;
\item \textbf{Lens:} quotient, redundancy, overlap, or multiple
representatives;
\item \textbf{Nil:} filtration, termination, nilpotent action, or a finite
survivor sector.
\end{itemize}
Each use must declare its category and object type.

\begin{proposition}[Taxonomy does not force a realization]
A CLN label does not by itself select a connection, gauge group, Hamiltonian,
probability law, spacetime action, or quantum theory.
\end{proposition}

\begin{proof}
The labels describe broad structural profiles.  Inequivalent bundles,
operators, actions, and observable maps can share those profiles.  A selector
or additional theorem is therefore required.
\end{proof}

\subsection{Corrected B-series interpretation}

Program B0 no longer claims that CLN is exhaustive or that exactly three
responses are forced.  It gives a taxonomy and conditional minimal curvature
realizations \cite{Nero2026_MTT_B0_CircleLensNil}.  Programs B1--B3 then study
conditional gravity, gauge, and discrete-survivor realizations under explicit
additional hypotheses
\cite{Nero2026_MTT_B1_GravityEncoding,Nero2026_MTT_B2_GaugeEncoding,
Nero2026_MTT_B3_QuantizationEncoding}.

Literal $S^1\times\mathrm{Lens}_3\times\mathrm{Nil}_3$ is seven-dimensional
and is not the selected six-dimensional q79 compactification.  Literal
manifold nesting is not inferred from a rank flag.  Lens and Nil may instead
appear as parallel circle-bundle models, filtration labels, or operator
profiles.

\section{The Corrected Structural Program A0--D1}

\subsection{A-series: reduced description}

\begin{itemize}
\item \textbf{A0} supplies the typed reduction dictionary, descent theorem,
basin-local fixed-point control, and explicit admissibility margins.
\item \textbf{A1} develops coherent kinematics conditionally from chart
persistence and declared overlap data; it does not derive Lorentzian dynamics
from projection alone.
\item \textbf{A2} analyzes conditional computability and finite prediction
depth; it does not prove universal undecidability at every boundary.
\end{itemize}

\subsection{B-series: conditional encoding responses}

\begin{itemize}
\item \textbf{B0} gives the non-exhaustive CLN taxonomy.
\item \textbf{B1} studies loop-transport consistency and a conditional gravity
realization.
\item \textbf{B2} separates gauge redundancy, global sections, and conditional
Yang--Mills realization.
\item \textbf{B3} proves discrete survivor-filter statements and records the
additional assumptions needed for quantum reconstruction.
\item \textbf{B4} distinguishes compatibility, local rigidity, persistence,
and global uniqueness.  Selected finite Standard Model packets establish a
particular compatibility branch, not uniqueness across all realizations
\cite{Nero2026_MTT_B4_StructuralRigidity}.
\item \textbf{B5} treats saturation relative to a declared contract.  Extended
carriers, worldsheets, and dualities require explicit additional structures
\cite{Nero2026_MTT_B5_SaturatedEncodings}.
\end{itemize}

\subsection{C-series: typed realization}

Program C is the realization dictionary.  It distinguishes geometry, bundles,
connections, operators, projectors, and finite carriers.  It proves joint
projector and compact-resolvent results under their stated hypotheses and
proves that realization nonuniqueness limits prediction until a source law
selects one realization \cite{Nero2026_MTT_C_Realizations}.

\subsection{D1: interpretive hypothesis}

Program D1 now treats E8/E9 dark-sector language as a projection-first
hypothesis only.  Encoding availability does not determine a stress tensor,
equation of state, or acceleration.  A covariant source or modified field
equation and observational execution remain necessary
\cite{Nero2026_MTT_D1_DarkSector}.

\section{Proof Tiers for Relations to Established Physics}

\subsection{Tier definitions}

We use the following hierarchy:
\begin{center}
\begin{tabular}{|c|p{10cm}|}
\hline
Level & Meaning \\ \hline
L0 & Interpretive analogy, diagnostic reframing, or conjecture. \\ \hline
L1 & Typed representation or compatibility map. \\ \hline
L2 & Conditional reconstruction after target-compatible structures are
assumed. \\ \hline
L3 & Controlled calculation, profile replay, or certified result on a selected
finite branch. \\ \hline
L4 & Selected physical source theorem with dynamics, state, observables,
uncertainty, and held-out validation. \\ \hline
\end{tabular}
\end{center}

\begin{remark}
The same research program may contain results at several levels.  A finite L3
packet does not automatically promote the surrounding continuum theory to L4.
\end{remark}

\subsection{Framework status map}

The following table records the strongest defensible corpus-level
interpretation, not a claim that every paper in a row is equally strong.

\begin{center}
\begin{tabular}{|
  >{\raggedright\arraybackslash}p{2.6cm}|
  >{\raggedright\arraybackslash}p{4.6cm}|
  >{\raggedright\arraybackslash}p{5.0cm}|}
\hline
Framework & Current contribution & Remaining physical obligation \\ \hline
Quantum mechanics & Complex Hilbert, operator, and selected finite recorder
structures can be represented or conditionally reconstructed. & A general
same-source Born theorem for all declared apparatus contexts remains open. \\
\hline
QFT and AQFT & Local nets, BV/BRST data, and perturbative structures exist
under imported locality, state, and renormalization hypotheses. & A selected
nonperturbative physical completion, state, RG matching, and observables remain
open. \\
\hline
General relativity & Typed geometric and source/descent relations can recover
Einstein-form dynamics conditionally. & Projection alone does not derive the
Einstein action, matter source, or all GR phenomenology. \\
\hline
Standard Model & Selected finite representation, gauge-group, anomaly, and
profile/parity packets establish substantial compatibility. & Zero-primitive
electroweak normalization and held-out no-knob precision equivalence remain
open. \\
\hline
Strings and flux geometry & q79 and Fu--Yau-type data provide a selected
candidate branch and exact finite constraints. & Physical visible-hidden HYM
endpoints, continuum naturality, action transfer, and the complete worldsheet
contract remain open. \\
\hline
Quantum gravity & Several fixed-point, perturbative, Euclidean, BRST, and
projection-first constructions are conditional or controlled. & UV-complete,
positive, causal, nonperturbative physical dynamics are not established by the
roadmap. \\
\hline
Dark sector & Availability profiles and their non-identifiability are typed. &
A selected covariant source, cosmological solutions, perturbations, nonlinear
tests, and likelihood remain open. \\
\hline
\end{tabular}
\end{center}

\subsection{Why ``MTT to X'' is not one claim}

A paper titled ``MTT to X'' may do one or more of the following:
\begin{itemize}
\item embed the notation of $X$ into an MTT carrier;
\item show compatibility with an $X$-like constraint;
\item reconstruct $X$ after importing its action or state space;
\item calculate a selected finite profile;
\item derive a source law and predict held-out observables.
\end{itemize}
Only the last item is a full physical derivation.  Titles do not determine
tiers; hypotheses and exit certificates do.

\section{Overlaps and Shadow Bridges}

\subsection{A lawful overlap}

An overlap between encodings $\mathfrak C_1$ and $\mathfrak C_2$ requires:
\begin{enumerate}[label=\arabic*.]
\item a common source domain or an explicit comparison span;
\item typed maps into both encodings;
\item a commuting relation on the declared domain;
\item preservation of the relevant dynamics or constraints;
\item an observable comparison and an error budget.
\end{enumerate}

For example, if $J_i:\mathcal U\to\mathcal E_i$ and the encoded evolutions are
$F_i$, a dynamics-preserving overlap requires a source evolution $F_{\mathcal
U}$ such that
\[
J_iF_{\mathcal U}=F_iJ_i
\]
exactly or within a declared controlled error.  Similar words in two papers
do not supply this diagram.

\subsection{Duality and equivalence}

A duality requires an invertible comparison on its declared physical quotient
and preservation of dynamics and observables.  A common rank count, shared
circle, matching spectrum fragment, or common low-energy limit is weaker.

\subsection{Shadow-bridge status}

A shadow bridge can be useful at L1 or L2 when it identifies a common
constraint or diagnostic.  It must be labeled accordingly.  It becomes a
physical equivalence only after its source, dynamics, state, observables, and
inverse properties are proved.

\section{Fixed Points, Actions, and Physical Promotion}

\subsection{Fixed points}

The corrected fixed-point sequence distinguishes:
\begin{itemize}
\item a fixed point of an auxiliary stabilization map;
\item a fixed point of a time-$\tau$ map;
\item a stationary solution of a physical evolution;
\item an attracting state under a Lyapunov or contractive law.
\end{itemize}
These are not equivalent without explicit hypotheses.  Banach arguments
require complete invariant basins; Schauder and Darbo give existence under
different compactness or condensing conditions; uniqueness needs coercivity,
strong monotonicity, or another declared mechanism.

\subsection{Action and source promotion}

To promote a geometric or operator realization physically, supply:
\begin{enumerate}[label=\textbf{P\arabic*.}]
\item a selected action, Hamiltonian, generator, or covariant field equation;
\item the state space, measure, or positive functional;
\item gauge constraints and anomaly conditions;
\item the observable map;
\item parameter and normalization provenance;
\item well-posedness and stability;
\item approximation and numerical error certificates;
\item comparison data separated into construction, calibration, and held-out
sets.
\end{enumerate}

\begin{proposition}[Realization compatibility is not prediction]
If two inequivalent realizations satisfy the same stated premises but give
different values for an observable, those premises do not predict that
observable.
\end{proposition}

\begin{proof}
Prediction requires the same value in every model of the premises, or a
selector that removes the alternatives.  Two admissible countermodels with
different values disprove that implication.
\end{proof}

\section{Current Frontier Snapshot}

\subsection{Closed control rows and open physical rows}

As of July 2026, the curated ledger records important closed control results,
including canonical-descent and GR-control rows, alongside open physical
source rows.  The principal open dependency chain includes:
\begin{itemize}
\item concrete eta9 meridian/period and flat Deligne values;
\item selected visible-hidden Hull--Strominger endpoints;
\item continuum q79 geometry-to-operator naturality;
\item selected continuum operator execution;
\item upper action and automorphism transfer;
\item the general Born source theorem;
\item nonperturbative QFT/BV completion;
\item the complete q79 heterotic worldsheet contract;
\item zero-primitive electroweak normalization;
\item no-knob Standard Model values and precision equivalence.
\end{itemize}

This list is a dated snapshot.  The live A/B rows, source hashes, and exit
certificates supersede it when they change.

\subsection{What finite success currently means}

Selected finite packets can establish exact arithmetic, representation,
anomaly, group, or profile statements on their declared branch.  They do not
by themselves establish:
\begin{itemize}
\item a physical continuum compactification;
\item an HYM source connection;
\item a selected upper action;
\item a general probability law;
\item a nonperturbative QFT;
\item held-out Standard Model precision.
\end{itemize}

\section{How to Read and Extend the Corpus}

\subsection{Paper audit checklist}

For every theorem or numerical claim, record:
\begin{enumerate}[label=\arabic*.]
\item the selected source path, version, and hash;
\item the exact premises and imported target structures;
\item domains and codomains of every map;
\item operator domains and commutation assumptions;
\item the proof tier L0--L4;
\item parameter provenance and units;
\item verifier and artifact paths;
\item open blockers and exit certificates;
\item whether data are construction inputs, calibration, replay, or held out.
\end{enumerate}

\subsection{Extension rule}

New work should change a declared frontier truth value, discharge an exit
certificate, prove a new conditional theorem, construct a missing typed map,
or provide an independently verified calculation.  Renaming an existing open
packet, repeating a compatibility result, or replacing an old benchmark with
another fitted replay is not frontier progress.

\subsection{Why the static paper index was removed}

Version 11 contained a long paper-by-paper table with statuses such as
``closed'' and ``mostly closed.''  That table became stale as successor
papers, corrigenda, and calculations changed.  Version 12 replaces it with:
\begin{itemize}
\item stable conceptual classes in this paper;
\item per-paper revision audits in the canonical paper repository;
\item current A/B rows and source health in the research kernel;
\item result references to the curated reproducibility repository where
applicable.
\end{itemize}
This separates exposition from live authority.

\section{Version 12 Revision Record}

Version 12 makes the following substantive corrections:
\begin{enumerate}[label=\arabic*.]
\item ordinary noninjectivity no longer implies absence of a right inverse;
\item descent, representative sections, exact decoding, merger, and stable
continuation are typed separately;
\item chart failure is expressed through named margins and means loss of
control for that chart, not universal termination of description;
\item the canonical physical convention is
$\pi:M_{10}\to Y_4$ with six-dimensional fiber $X_6$;
\item vertical triplet operators, line bundles, and coordinate factors are
not conflated;
\item the shared circle is phase/holonomy data, not time or an extra
coordinate;
\item CLN is a non-exhaustive taxonomy and does not force unique physical
responses;
\item the ten-dimensional statement is replaced by the conditional
transverse-curvature Proposition~\ref{prop:fourplussix};
\item QM, QFT, GR, SM, strings, QG, AQFT, and dark-sector relations are
classified by proof tier and remaining source obligations;
\item the stale static closure index is replaced by a stable reading guide and
live-ledger authority rule.
\end{enumerate}

\section{Conclusion}

The coherent core of MTT is best understood as a typed research program about
conditional reduced description.  Its strongest general result is not that
every familiar theory follows automatically, but that projection, descent,
recovery, coherent continuation, and physical selection can be separated and
tested with explicit contracts.

The canonical physical realization uses a four-dimensional Lorentzian base
and a six-dimensional internal fiber when that specialization is chosen.
Circle-Lens-Nil organizes recurring structures but neither exhausts all
obstructions nor selects gravity, gauge theory, or quantum mechanics by name.
Framework reconstructions are valuable at their proved tier; they become full
physical derivations only when the selected source, dynamics, state,
observables, uncertainty, and held-out tests are supplied.

This roadmap therefore offers a firmer unity than the old closure language:
every paper can be located by its types, dependencies, proof tier, and exit
certificate.  That structure preserves genuine achievements while keeping the
remaining frontier visible.

\bibliographystyle{plain}
\bibliography{main}

\end{document}
